\section{Introduction}
\label{sec:introduction}

Executable algebra in a proof assistant needs concrete algorithms and exact
mathematical interfaces. Matrix normal forms make the tension especially
visible: elimination uses local row and column operations, but the result must
record a canonical matrix and a transformation whose invertibility is trusted
by the kernel.

This project develops \HNF{} and \SNF{} in \Lean{} against \mathlib{}
\citep{deMouraUllrich2021Lean4,mathlib2020}. Version 1.2.2 of the accompanying
software artifact fixes the five public interfaces below
\citep{ji2026normalforms}.

\begin{itemize}
  \item A strong certificate stores a chosen inverse and proofs of both
    multiplication identities. HNF and SNF results carry these certificates
    end to end, including elementary operations, Bezout blocks, recursive
    lifts, and composition.
  \item Normal-form semantics never hide an enumeration of an arbitrary finite
    type. Public callers may choose an indexing explicitly, supply
    equivalences, request increasing order, or use a convenience arbitrary
    indexing that is retained in the result.
  \item Executable branching is isolated behind
    \code{ComputableEuclideanOps}. The matrix kernels retain ordinary ring
    arithmetic, while equality, division with remainder, normalization,
    extended gcd, and the termination measure have constructive
    implementations and semantic laws.
  \item External generators return only raw data. Following the skeptical
    integration pattern for computer-algebra output
    \citep{harrisonThery1998}, a pure checker binds the
    certificate to a caller-supplied matrix, and strict generated modules reduce
    each obligation with \code{decide_cbv} before obtaining a strong result
    through the checker soundness theorem.
  \item Intrinsic Smith data and finite-free presentations omit execution
    traces. Determinantal ideals characterize rank and associate factors, and a
    presentation matrix always maps relations to generators by columns. The
    Fitting-ideal convention follows the minors of a presentation matrix
    \citep{stacksFittingIdeals}.
\end{itemize}

The \code{Fin}-indexed HNF and SNF kernels are complete with correctness and
canonicality theorems. HNF uniqueness is deliberately stated only for one
fixed indexing; changing column order yields a constructive left equivalence,
not an unjustified equality. Smith uniqueness is strong enough to prove that
the canonical \code{Fin} Smith matrix and its invariant factors are independent
of indexing. \code{SmithData} records the exact nonzero prefix and full
diagonal, while \code{SmithSignature} records all nonzero associate factors.
Determinantal ideals prove that the signature is invariant under certified
transforms and reindexing. A canonical wrapper compares it with mathlib's PID
Smith witnesses. \code{FiniteFreePresentation} fixes columns as relations,
defines the cokernel and Fitting ideals, and derives a Smith decomposition.
Both integer and rational-polynomial matrices execute through deterministic
\code{\#eval} and native baselines; the polynomial implementation includes
direct finite-support arithmetic, long division, monic normalization, and
well-founded extended gcd.

The optional FLINT 3.6.0 adapter offers two untrusted routes. A process worker
and a separately linked Lean FFI call the same C computation and strict decimal
protocol; the Python standard-library wrapper emits
\code{normalforms.cert/v1}. FLINT, GMP, and MPFR are pinned and dynamically
linked, but none is a theorem dependency. Both routes pass through the same
pure checker. Lean independently checks the expected input, transformation
equation, both inverse directions, and normal-form predicate.

Fixed-seed differential regressions compare normalized invariant factors, not
the nonunique transformation matrices. A recorded performance profile
separates internal Lean generation, process and FFI generation, native
checking, strict kernel checking, and complete generation/import costs. Its
raw observations and hardware fingerprint make clear that generator or native
success is experimental evidence; only recompiling a strict generated module
constructs a theorem.

Divas\'on and Thiemann formalize Smith normal form in Isabelle/HOL, and Cano
et al. formalize linear algebra over elementary divisor rings in Coq
\citep{divasonThiemann2022snf,cano2016coq}. Ma, Wang, and Wen give a Lean
framework whose bridge-derived Smith and rational-canonical results are
intentionally noncomputable and do not establish normalized uniqueness or an
executable reduction trajectory \citep[Sections 3.4, 5.1, and
6.2]{maWangWen2026}. The contribution here is an executable Lean/\mathlib
architecture with explicit inverses, precise indexing semantics, canonicality,
a PID bridge, and a certifying external-computation path.

The version 1.2.2 release freezes the facade and presentation direction. Exact
type tests, synchronized metadata checks, separate executable-core and
certificate axiom audits, and a pinned artifact profile make that boundary
machine-checkable.
